\chapter{ซอร์สโค้ดฉบับสมบูรณ์ของระบบ SoilpH-HT AI}
\label{app:code}

ภาคผนวกนี้รวบรวมซอร์สโค้ดภาษา Dart ฉบับสมบูรณ์ของโมดูลแกนหลักในระบบ \textbf{SoilpH-HT AI} ซึ่งได้รับการออกแบบและคอมไพล์เพื่อใช้งานบนอุปกรณ์สมาร์ตโฟนจริงผ่านพอร์ต USB Type-C OTG

\section{โมดูลบริการสื่อสารผ่านพอร์ตฮาร์ดแวร์จริง (usb\_ph\_service.dart)}

\begin{codebox}[title={ไฟล์ lib/data/services/usb\_ph\_service.dart (คัดย่อส่วนสำคัญ)}]
\begin{lstlisting}[language=Java]
import 'dart:async';
import 'dart:typed_data';
import 'package:flutter/foundation.dart';
import 'package:usb_serial/usb_serial.dart';
import '../models/modbus_parser.dart';
import '../models/soil_ph_reading.dart';

enum SensorConnectionStatus { disconnected, connecting, connected, error }

class UsbPhService extends ChangeNotifier {
  UsbPort? _port;
  StreamSubscription<Uint8List>? _transactionSub;
  Timer? _pollingTimer;
  Timer? _watchdogTimer;

  int _baudRate = 4800; // ค่าเริ่มต้นมาตรฐาน Modbus
  int _txCount = 0;
  int _rxByteCount = 0;
  String _lastRxHex = '---';

  Future<bool> connectToDevice(UsbDevice device) async {
    _port = await device.create();
    if (_port == null) return false;

    final openSuccess = await _port!.open();
    if (!openSuccess) return false;

    await _port!.setDTR(true);
    await _port!.setRTS(false); // ปลด RE เป็น Low เพื่อเปิดรับสัญญาณ
    await _port!.setPortParameters(
      _baudRate,
      UsbPort.DATABITS_8,
      UsbPort.STOPBITS_1,
      UsbPort.PARITY_NONE,
    );
    await _port!.setFlowControl(UsbPort.FLOW_CONTROL_OFF);

    _transactionSub = _port!.inputStream!.listen(_onDataReceived);
    _startPolling();
    return true;
  }

  void _startPolling() {
    _pollingTimer?.cancel();
    _pollingTimer = Timer.periodic(const Duration(milliseconds: 1500), (_) {
      if (_port == null) return;
      final requestFrame = ModbusParser.buildReadRequest(
        slaveId: 0x01,
        startRegister: 0x0000,
        registerCount: 4, // อ่านความชื้น อุณหภูมิ EC และ pH
      );
      _port!.write(requestFrame);
      _txCount++;
      _resetWatchdog();
      notifyListeners();
    });
  }

  void _onDataReceived(Uint8List chunk) {
    _rxByteCount += chunk.length;
    _lastRxHex = chunk.map((b) => b.toRadixString(16).padLeft(2, '0')).join(' ');
    
    final reading = ModbusParser.parsePhResponse(chunk);
    if (reading != null) {
      _watchdogTimer?.cancel();
      notifyListeners();
    }
  }

  void _resetWatchdog() {
    _watchdogTimer?.cancel();
    _watchdogTimer = Timer(const Duration(seconds: 3), () {
      _baudRate = (_baudRate == 4800) ? 9600 : 4800;
    });
  }
}
\end{lstlisting}
\end{codebox}

\section{โมดูลถอดรหัสโปรโตคอลและตรวจสอบความถูกต้อง (modbus\_parser.dart)}

\begin{codebox}[title={ไฟล์ lib/data/models/modbus\_parser.dart}]
\begin{lstlisting}[language=Java]
import 'dart:typed_data';
import 'soil_ph_reading.dart';

class ModbusParser {
  ModbusParser._();

  /// คำนวณรหัสตรวจสอบความถูกต้อง CRC-16 Modbus (พหุนาม 0xA001)
  static int calculateCrc16(List<int> data) {
    int crc = 0xFFFF;
    for (int byte in data) {
      crc ^= byte & 0xFF;
      for (int i = 0; i < 8; i++) {
        if ((crc & 0x0001) != 0) {
          crc = (crc >> 1) ^ 0xA001;
        } else {
          crc >>= 1;
        }
      }
    }
    return crc & 0xFFFF;
  }

  /// สร้างเฟรมคำสั่งอ่าน Holding Registers
  static Uint8List buildReadRequest({int slaveId = 1, int startRegister = 0, int registerCount = 4}) {
    final buffer = Uint8List(8);
    final bd = ByteData.sublistView(buffer);
    buffer[0] = slaveId;
    buffer[1] = 0x03;
    bd.setUint16(2, startRegister, Endian.big);
    bd.setUint16(4, registerCount, Endian.big);
    final crc = calculateCrc16(buffer.sublist(0, 6));
    bd.setUint16(6, crc, Endian.little);
    return buffer;
  }

  /// ถอดรหัสเฟรมตอบกลับ Modbus RTU สำหรับ pH อุณหภูมิ และความชื้น
  static SoilPhReading? parsePhResponse(Uint8List frame) {
    if (frame.length < 13) return null;
    final payload = frame.sublist(0, frame.length - 2);
    final calculatedCrc = calculateCrc16(payload);
    final receivedCrc = frame[frame.length - 2] | (frame[frame.length - 1] << 8);
    if (calculatedCrc != receivedCrc) return null;

    final bd = ByteData.sublistView(frame);
    final moisture = bd.getUint16(3, Endian.big) / 10.0;
    final temperature = bd.getInt16(5, Endian.big) / 10.0;
    final ec = bd.getUint16(7, Endian.big).toDouble();
    final rawPh = bd.getUint16(9, Endian.big) / 10.0;

    return SoilPhReading(
      ph: rawPh,
      temperature: temperature,
      moisture: moisture,
      ec: ec,
      timestamp: DateTime.now(),
    );
  }
}
\end{lstlisting}
\end{codebox}

\section{โมดูลเอนจินการคำนวณและชดเชยค่าอิงฟิสิกส์ (deep\_learning\_calibrator.dart)}

\begin{codebox}[title={ไฟล์ lib/domain/calibration/deep\_learning\_calibrator.dart}]
\begin{lstlisting}[language=Java]
import 'dart:math' as math;

class DeepLearningCalibrator {
  DeepLearningCalibrator._();

  static const double tRef = 25.0; // อุณหภูมิอ้างอิงมาตรฐาน 25 องศาเซลเซียส

  /// การปรับเทียบความคลาดเคลื่อนสองขั้นตอน (Two-Stage Decoupling)
  static Map<String, double> calibratePh({
    required double rawPh,
    required double rawTemp,
    required double rawMoisture,
    double rawEc = 0.0,
  }) {
    final deltaT = rawTemp - tRef;

    // ขั้นตอนที่ 1 ฟิสิกส์หลักการแรก (First-Principles Physics)
    final nernstRatio = (273.15 + tRef) / (273.15 + rawTemp);
    final deltaPhThermal = (rawPh - 7.0) * (nernstRatio - 1.0);
    final deltaPhWater = 0.017 * (25.0 - rawTemp);
    final deltaPhT = deltaPhThermal + deltaPhWater;

    // ชดเชยอิมพีแดนซ์รอยต่อในดินแห้ง
    final deltaPhM = 0.35 * math.exp(-0.08 * rawMoisture.clamp(0.0, 100.0));

    // ขั้นตอนที่ 2 โครงข่ายประสาทเทียม PINN ชดเชยปฏิสัมพันธ์ร่วมไม่เป็นเชิงเส้น
    final normT = deltaT / 20.0;
    final normM = (rawMoisture - 50.0) / 50.0;
    final h1 = _swish(0.45 * normT - 0.38 * normM + 0.12);
    final h2 = _gelu(0.28 * h1 + 0.35 * normT * normM - 0.05);
    final deltaPhPinn = 0.08 * h2;

    // ป้องกันค่าหลุดขอบเขตด้วยการ Clamp
    final calibratedPh = (rawPh + deltaPhT + deltaPhM + deltaPhPinn).clamp(3.0, 10.5);

    return {
      'rawPh': rawPh,
      'calibratedPh': calibratedPh,
      'deltaPh': calibratedPh - rawPh,
      'deltaPhT': deltaPhT,
      'deltaPhM': deltaPhM,
      'deltaPhPinn': deltaPhPinn,
      'temperature': rawTemp,
      'moisture': rawMoisture,
    };
  }

  static double _swish(double x) => x / (1.0 + math.exp(-x));
  static double _gelu(double x) => 0.5 * x * (1.0 + math.tanh(0.79788 * (x + 0.044715 * x * x * x)));
}
\end{lstlisting}
\end{codebox}
